We interpret the problem in the standard number-theoretic way:
\emph{$k,m,n$ are nonnegative integers}, with $k\ge 2$ and $m\ge n+k$.
For such $(k,m,n)$ define
\[
R(k,m,n):=\frac{\prod_{i=1}^k (m+i)}{\prod_{i=1}^k (n+i)}.
\]
Question: Is it true that for every integer $N\ge 2$ there exist integers $k\ge 2$ and $m\ge n+k\ge 0$ such that
$R(k,m,n)=N$?

\emph{Ambiguity note.} If one allows negative $n$ or $m$, then denominators can vanish or signs can appear; the problem statement
suggests the intended domain is nonnegative integers with $m\ge n+k$.
